\chapter*{Introduction Générale}
\addcontentsline{toc}{chapter}{Introduction Générale}
\markboth{INTRODUCTION GÉNÉRALE}{INTRODUCTION GÉNÉRALE}

Aujourd'hui, le numérique transforme profondément tous les secteurs d'activité. Les cafés
et restaurants ne font pas exception : digitaliser l'expérience client est devenu un avantage
concurrentiel réel pour les petits établissements qui cherchent à moderniser leur service
sans engager des coûts importants.

C'est dans ce contexte que s'inscrit notre projet de fin d'études, intitulé \textbf{Coffee Time}.
Il s'agit d'une application web et mobile multiservices destinée à digitaliser l'expérience
client dans un coffee shop. Son principe est simple : chaque table est équipée d'un QR code
unique. En le scannant, le client accède directement à l'application depuis son smartphone,
sans inscription ni connexion, et une session liée à sa table est créée automatiquement.

La plateforme regroupe quatre modules complémentaires : un système de commande en
ligne, un module de divertissement interactif, un module d'actualités via flux RSS, et un
tableau de bord administrateur destiné au personnel du café.

Sur le plan technique, la solution repose sur React.js pour le frontend web, React Native
pour le mobile, Node.js et Express pour le backend, et PostgreSQL comme base de données
relationnelle. La communication entre les différentes couches s'effectue via une API REST.
Le projet a été conduit selon la méthodologie Agile Kanban, organisée en quatre phases de
développement successives.

Le présent rapport est structuré comme suit :

\begin{itemize}
    \item \textbf{Chapitre 1 --- Contexte général et étude de l'existant :} présentation de
    l'organisme d'accueil, analyse des solutions existantes sur le marché et justification
    de la solution proposée.

    \item \textbf{Chapitre 2 --- Analyse et spécification des besoins :} identification des acteurs,
    besoins fonctionnels et non fonctionnels, backlog des tâches, organisation en phases,
    architecture du système et environnement logiciel.

    \item \textbf{Chapitre 3 --- Réalisation Phase 1 et Phase 2 :} réalisation des deux premières
    phases. La Phase~1 couvre l'authentification et la gestion des accès. La Phase~2 traite
    le module de commande côté client, incluant la session QR code, le menu, le panier
    et le suivi de commande en temps réel.

    \item \textbf{Chapitre 4 --- Réalisation Phase 3 et Phase 4 :} réalisation des deux dernières
    phases. La Phase~3 porte sur le tableau de bord d'administration et la gestion du
    personnel. La Phase~4 couvre les services additionnels : divertissement, actualités RSS,
    programme de fidélité et chatbot intelligent basé sur l'intelligence artificielle.
\end{itemize}
